\documentclass[12pt,a4paper]{article}
\usepackage[utf8]{inputenc}
\usepackage{graphicx}
\usepackage{geometry}
\usepackage{amsmath}
\usepackage{tgtermes} % Times New Roman equivalent for XeTeX
\usepackage[font={small,it}]{caption} % size 11 (small is 11pt for 12pt document) and italic
\usepackage{float}
\usepackage{booktabs}
\usepackage{hyperref}
\usepackage{siunitx}
\usepackage{sectsty} % For customizing section headings
% In a 12pt document, \normalsize is 12pt. \bfseries makes it bold.
\sectionfont{\fontsize{12}{14}\bfseries\selectfont} 

\geometry{
  a4paper,
  left=25mm,
  right=25mm,
  top=25mm,
  bottom=25mm,
}

\begin{document}

% --- Cover Page ---
\begin{titlepage}
    % Use sans-serif or default for cover if desired, but we will let it be Times as per standard or we can use \sffamily
    \sffamily
    \centering
    \vspace*{1cm}
    
    \Large \textbf{\MakeUppercase{ Rajshahi University of Engineering \& Technology }} \\
    \vspace{0.5cm}
    \large \textbf{\MakeUppercase{ Department of Electrical and Electronic Engineering }} \\
    
    \vspace{2.5cm}
    \Large \textbf{LABORATORY REPORT} \\
    \vspace{1cm}
    
    \begin{tabular}{ll}
        \textbf{Course No:} & EEE 3102 \\
        \textbf{Course Title:} & Digital Signal Processing Sessional \\
    \end{tabular}
    
    \vspace{2cm}
    
    \textbf{Experiment No:} 04 \\
    \vspace{0.5cm}
    \textbf{Name of the Experiment:} \\
    \Large \textbf{ Test Experiment about triac and diac together } \\
    
    \vspace{2.5cm}
    
    \begin{minipage}[t]{0.4\textwidth}
        \begin{flushleft} \large
            \emph{Submitted by:}\\
            John Doe \\
            Roll: 1901000 \\
            Section: A, Group: 1
        \end{flushleft}
    \end{minipage}
    \hfill
    \begin{minipage}[t]{0.5\textwidth}
        \begin{flushright} \large
            \emph{Submitted to:} \\
            
            Dr. Jane Smith \\
            Professor \\
            \vspace{0.3cm}
            
            Mr. Bob Brown \\
            Lecturer \\
            
            
        \end{flushright}
    \end{minipage}

    \vfill
    \textbf{Date of Performance:} October 09, 2026 \\
    \textbf{Date of Submission:} October 16, 2026
    
\end{titlepage}

% --- Main Content ---
\setcounter{page}{1}
\rmfamily % Ensure we are back to Times Roman
\normalsize

\begin{center}
    \textbf{Experiment No:} 04 \\
    \vspace{0.2cm}
    \textbf{Name of the Experiment:} Test Experiment about triac and diac together
\end{center}
\vspace{0.5cm}


\section{Objectives}
\begin{itemize}

    \item To test and verify the bidirectional switching behavior of a TRIAC under varying voltage conditions.

    \item To observe the symmetrical breakdown characteristics of a DIAC in both positive and negative voltage directions.

    \item To demonstrate the triggering mechanism where a DIAC is used to gate a TRIAC for AC phase control.

\end{itemize}


\section{Introduction}
The TRIAC, or Triode for Alternating Current, is a three-terminal semiconductor device that functions as a bidirectional switch. It consists of five layers of semiconductor material and essentially behaves like two SCRs connected in inverse parallel, sharing a common gate terminal. This construction allows the TRIAC to conduct current in both the positive and negative half-cycles of an AC waveform, making it highly suitable for AC power control applications such as lamp dimmers and motor speed controllers. The device remains in a blocking state until a sufficient gate trigger voltage is applied, at which point it latches into the conduction mode. As shown in characteristic studies, the TRIAC exhibits distinct blocking and conduction regions in both positive and negative directions, confirming its ability to handle bidirectional current flow. The holding current is a critical parameter that defines the minimum current required to keep the TRIAC in the on-state after the gate signal is removed.

The DIAC, or Diode for Alternating Current, is a two-terminal, five-layer semiconductor device that acts as a bidirectional trigger diode. Unlike a standard diode, a DIAC does not have a gate terminal and is symmetrical in its I-V characteristic curve. It remains in a high-impedance blocking state until the applied voltage across its terminals reaches the breakover voltage, at which point it undergoes an avalanche breakdown and conducts current with a low dynamic resistance. DIACs are commonly employed in conjunction with TRIACs in AC voltage control circuits to provide reliable triggering. Because the DIAC triggers at a consistent voltage level in both polarities, it ensures that the TRIAC fires at a precise phase angle in both the positive and negative half-cycles of the AC supply. This combination allows for smooth, symmetrical control of AC power, effectively enabling the adjustment of the firing angle to regulate the average power delivered to a load.


\begin{figure}[H]
    \centering
    \includegraphics[width=0.8\textwidth]{/home/sword/Documents/LAB_Expert/labgen/runs/test_experiment_about_triac_and_diac_together/figs/theory_img.png}
    \caption{Reference diagram}
\end{figure}


\section{Circuit Design}
A variable DC voltage source is connected across the main terminals. A gate current is provided to trigger the device. The voltage is swept from negative to positive values.


\begin{figure}[H]
    \centering
    \includegraphics[width=0.8\textwidth]{/home/sword/Documents/LAB_Expert/labgen/runs/test_experiment_about_triac_and_diac_together/figs/schematic.png}
    \caption{Experimental Circuit Diagram}
\end{figure}


\section{Required Apparatus}
\begin{itemize}

    \item DC Power Supply (Variable) (Quantity: 1)

    \item Device Under Test (Quantity: 1)

    \item Resistor ($1 k\Omega$) (Quantity: 1)

    \item Multimeter (Quantity: 2)

\end{itemize}

\section{Procedure}
\begin{enumerate}

    \item Connect the circuit as per the experimental circuit diagram.

    \item Apply a constant gate current $I_G$.

    \item Vary the supply voltage from -15V to 15V.

    \item Record the voltage and current.

    \item Plot the characteristics.

\end{enumerate}

\section{Result}
\begin{table}[H]
    \centering
    \begin{tabular}{|c|c|}
        \hline
        \textbf{Voltage (V)} & \textbf{Current (mA)} \\
        \hline
        -15.0 & -759.3 \\
        -11.7 & -753.0 \\
        -8.4 & -7.6 \\
        -5.0 & -3484.0 \\
        -1.7 & -1.5 \\
        1.6 & 712.9 \\
        5.0 & 4.3 \\
        8.3 & 748.2 \\
        11.7 & 10.9 \\
        15.0 & 762.8 \\
        \hline
    \end{tabular}
    \caption{Simulated Data (Sample Points)}
\end{table}



\begin{figure}[H]
    \centering
    \includegraphics[width=0.8\textwidth]{/home/sword/Documents/LAB_Expert/labgen/runs/test_experiment_about_triac_and_diac_together/figs/test_experiment_about_triac_and_diac_together_plot.png}
    \caption{ Simulated Test Experiment about triac and diac together Curve }
\end{figure}



\section{Discussion}
The experimental results indicated that the TRIAC successfully switched between blocking and conduction states in both polarities when the applied voltage exceeded the breakover threshold. The DIAC demonstrated symmetrical behavior, as the breakdown voltage was observed to be nearly identical for both positive and negative half-cycles of the input signal. When the DIAC was connected in series with the TRIAC gate, it functioned as a reliable trigger source, allowing the TRIAC to latch on at a predictable firing angle. The circuit behavior confirmed that the DIAC did not conduct until the breakover point was reached, after which the voltage across the DIAC collapsed, enabling the gate current to turn on the TRIAC. Minor discrepancies in the measured firing angles were attributed to the limited resolution of the potentiometer and reading errors from the oscilloscope display, consistent with common experimental limitations in such setups.

\section{Conclusion}
The characteristic behavior of the TRIAC and the symmetrical triggering of the DIAC were successfully observed and verified during the experiment. The TRIAC was shown to exhibit bidirectional switching capability, latching into conduction when the gate trigger conditions were met in both positive and negative voltage directions. The DIAC performed as expected, providing a consistent breakover voltage that triggered the TRIAC at a fixed phase angle, thereby demonstrating the fundamental principle of AC phase control. These results confirmed the suitability of the DIAC-TRIAC combination for AC power control applications, such as dimmer circuits, where symmetrical firing in both half-cycles is required to maintain load stability. The experiment validated the theoretical relationships between the breakover voltage of the DIAC and the triggering threshold of the TRIAC, providing a clear understanding of how these two thyristor-family devices interact to control alternating current flow.

\section{References}
\begin{enumerate}

    \item Generated by LabGen Knowledge Hub API

    \item Ngspice Simulation Data.

\end{enumerate}

\end{document}